\documentclass{report}
\usepackage{amsmath,amssymb}
\setlength{\parindent}{0mm}
\setlength{\parskip}{1em}
\begin{document}
\begin{center}
\rule{6in}{1pt} \
{\large D. Palitta \\
{\bf COMPUTATIONALLY ENHANCED PROJECTION METHODS FOR SYMMETRIC LYAPUNOV MATRIX EQUATIONS}}

Dipartimento di Matematica \\ Piazza di Porta San Donato 5 \\ 40126 \\ Bologna (Italy)
\\
{\tt davide.palitta3@unibo.it}\\
V. Simoncini\end{center}

In the numerical treatment of large-scale Lyapunov equations, projection
methods require solving a reduced Lyapunov problem to check convergence.
As the approximation space expands, this solution takes an increasing
portion of the overall computational effort. When data are symmetric, we
show that the Frobenius norm of the residual matrix can be computed at
significantly lower cost than with available methods, without explicitly
solving the reduced problem. For certain classes of problems, the new
residual norm expression combined with a memory-reducing device make
classical Krylov strategies competitive with respect to more recent
projection methods. In this talk, we present several numerical
experiments that illustrate the effectiveness of the new implementation
for standard and extended Krylov subspace methods.


\end{document}
